\documentclass[10pt,a4paper]{amsart}
\usepackage{amsmath,amssymb,amsthm,graphicx}
% For commutative diagrams we use Xy-pic
\usepackage[all]{xy}

% We use 2cell for 2-commutative diagrams.
\xyoption{2cell}
% Theorem environments.
%
\theoremstyle{plain}
\newtheorem{theorem}[subsection]{Theorem}
\newtheorem{proposition}[subsection]{Proposition}
\newtheorem{lemma}[subsection]{Lemma}

\theoremstyle{definition}
\newtheorem{definition}[subsection]{Definition}
\newtheorem{example}[subsection]{Example}
\newtheorem{exercise}[subsection]{Exercise}
\newtheorem{situation}[subsection]{Situation}

\theoremstyle{remark}
\newtheorem{remark}[subsection]{Remark}
\newtheorem{remarks}[subsection]{Remarks}

\numberwithin{equation}{subsection}

% Macros
%
\def\lim{\mathop{\mathrm{lim}}\nolimits}
\def\colim{\mathop{\mathrm{colim}}\nolimits}
\def\Spec{\mathop{\mathrm{Spec}}}
\def\Hom{\mathop{\mathrm{Hom}}\nolimits}
\def\Ext{\mathop{\mathrm{Ext}}\nolimits}
\def\SheafHom{\mathop{\mathcal{H}\!\mathit{om}}\nolimits}
\def\SheafExt{\mathop{\mathcal{E}\!\mathit{xt}}\nolimits}
\def\Sch{\mathit{Sch}}
\def\Mor{\mathop{\mathrm{Mor}}\nolimits}
\def\Ob{\mathop{\mathrm{Ob}}\nolimits}
\def\Sh{\mathop{\mathit{Sh}}\nolimits}
\def\NL{\mathop{N\!L}\nolimits}
\def\CH{\mathop{\mathrm{CH}}\nolimits}
\def\proetale{{pro\text{-}\acute{e}tale}}
\def\etale{{\acute{e}tale}}
\def\QCoh{\mathit{QCoh}}
\def\Ker{\mathop{\mathrm{Ker}}}
\def\Im{\mathop{\mathrm{Im}}}
\def\Coker{\mathop{\mathrm{Coker}}}
\def\Coim{\mathop{\mathrm{Coim}}}

% Boxtimes
%
\DeclareMathSymbol{\boxtimes}{\mathbin}{AMSa}{"02}

%
% Macros for moduli stacks/spaces
%
\def\QCohstack{\mathcal{QC}\!\mathit{oh}}
\def\Cohstack{\mathcal{C}\!\mathit{oh}}
\def\Spacesstack{\mathcal{S}\!\mathit{paces}}
\def\Quotfunctor{\mathrm{Quot}}
\def\Hilbfunctor{\mathrm{Hilb}}
\def\Curvesstack{\mathcal{C}\!\mathit{urves}}
\def\Polarizedstack{\mathcal{P}\!\mathit{olarized}}
\def\Complexesstack{\mathcal{C}\!\mathit{omplexes}}
% \Pic is the operator that assigns to X its picard group, usage \Pic(X)
% \Picardstack_{X/B} denotes the Picard stack of X over B
% \Picardfunctor_{X/B} denotes the Picard functor of X over B
\def\Pic{\mathop{\mathrm{Pic}}\nolimits}
\def\Picardstack{\mathcal{P}\!\mathit{ic}}
\def\Picardfunctor{\mathrm{Pic}}
\def\Deformationcategory{\mathcal{D}\!\mathit{ef}}
\usepackage{cleveref}
\crefname{originalreference}{original reference}{original references}
\Crefname{originalreference}{Original reference}{Original references}
\makeatletter
\newlabel{lemma-proper-base-change-stalk}{{0DDF}{source}{}{sourcepointer-1}{}}
\newlabel{lemma-proper-base-change-stalk@cref}{{[originalreference][0][]0DDF}{[1][0][]source}{}{}{}}
\makeatother
\crefname{originalreference}{original reference}{original references}
\Crefname{originalreference}{Original reference}{Original references}
\makeatletter
\newlabel{lemma-pull-out-constant}{{0F0E}{source}{}{sourcepointer-2}{}}
\newlabel{lemma-pull-out-constant@cref}{{[originalreference][0][]0F0E}{[1][0][]source}{}{}{}}
\makeatother
\crefname{originallemma}{original lemma}{original lemmas}
\Crefname{originallemma}{Original lemma}{Original lemmas}
\makeatletter
\newlabel{sites-cohomology-lemma-pullback-tensor-product}{{07A4}{source}{}{sourcepointer-3}{}}
\newlabel{sites-cohomology-lemma-pullback-tensor-product@cref}{{[originallemma][0][]07A4}{[1][0][]source}{}{}{}}
\makeatother
\crefname{originalsection}{original section}{original sections}
\Crefname{originalsection}{Original section}{Original sections}
\makeatletter
\newlabel{sites-cohomology-section-projection-formula}{{0943}{source}{}{sourcepointer-4}{}}
\newlabel{sites-cohomology-section-projection-formula@cref}{{[originalsection][0][]0943}{[1][0][]source}{}{}{}}
\makeatother
\title{The projection formula for a proper morphism with torsion coefficients over the ring of integers}
\author{The Stacks Project Authors}
\date{stacks-project:etale-cohomology.tex:local-snapshot-20261011; GNU FDL 1.2 or later}
\begin{document}
\maketitle
\noindent\textit{Editorial note.} All delivered mathematical text is the author's own, copied from the pinned Stacks Project chapter file etale-cohomology.tex; no statement or proof step was rewritten, reordered or removed.\par\medskip
\noindent\textit{Editorial note.} The counted proof is the author's complete argument: the proof is written out in full: the reduction to stalks with the commutation of the derived tensor product with pullback, the vanishing range from global dimension one, the torsion computation showing the coefficient sheaves stay torsion, and the application of the pull-out lemma.\par\medskip
\noindent\textit{Editorial note.} The label lemma-projection-formula-proper names the counted statement, and the proof environment immediately following it in the source is delivered verbatim as its proof.\par\medskip
\noindent\textit{Editorial note.} The counted statement is carried by the author's own lemma environment in the chapter source.\par\medskip
\noindent\textit{Editorial note.} The statement and proof are copied from the pinned chapter file; the four cited inputs are kept as external pointers with their canonical Stacks tags.\par\medskip
\noindent\textit{Editorial note.} The preamble is the chapter's own native preamble excerpt already used for this source: the Xy-pic setup (native preamble lines 14-18) and the native theorem declarations with the mathematical macros (native preamble lines 154-217).\par\medskip
\noindent\textit{Allowed standard prerequisites (editorial).}\begin{itemize}
\item The canonical projection map and its construction. Cohomology on Sites, the section on the projection formula, providing the canonical map Rf\_*E tensor\textasciicircum{}L K -> Rf\_*(E tensor\textasciicircum{}L f\textasciicircum{}\{-1\}K). Supplies the map whose being an isomorphism is proved.
\item Commutation of derived tensor products with pullback. Cohomology on Sites, the lemma identifying (E tensor\textasciicircum{}L f\textasciicircum{}\{-1\}K)\_x with E\_x tensor\textasciicircum{}L K\_y at geometric points. Supplies the stalk computation at the start of the proof.
\item Torsion coefficients and the derived tensor product with Q. Same chapter, the vanishing of the torsion of the tensor product and the pull-out lemma for a constant complex of abelian groups. Supplies that the cohomology sheaves of E tensor\textasciicircum{}L f\textasciicircum{}\{-1\}K are torsion and the vanishing range used with the pull-out lemma.
\end{itemize}\medskip
\ifdefined\hypertarget\hypertarget{sourcepointer-1}{}\fi
\noindent\textit{Source reference.} Original reference 0DDF: sources/stacks-etale-cohomology/pinned-original-source.tex, source lines 18734--18734. Prerequisite pointer invoked at proof line(s) 19041; it names a standard fact of the same chapter that the counted argument uses, and every step of the counted proof built on it is retained verbatim in the delivered body.\par\medskip
\ifdefined\hypertarget\hypertarget{sourcepointer-2}{}\fi
\noindent\textit{Source reference.} Original reference 0F0E: sources/stacks-etale-cohomology/pinned-original-source.tex, source lines 18938--18938. Prerequisite pointer invoked at proof line(s) 19050; it names a standard fact of the same chapter that the counted argument uses, and every step of the counted proof built on it is retained verbatim in the delivered body.\par\medskip
\ifdefined\hypertarget\hypertarget{sourcepointer-3}{}\fi
\noindent\textit{Source reference.} Original lemma 07A4: sources/stacks-chow/tags.txt, source lines 9294--9294. Prerequisite pointer invoked at proof line(s) 19018; it names a standard fact of a sibling chapter that the counted argument uses, and every step of the counted proof built on it is retained verbatim in the delivered body.\par\medskip
\ifdefined\hypertarget\hypertarget{sourcepointer-4}{}\fi
\noindent\textit{Source reference.} Original section 0943: sources/stacks-chow/tags.txt, source lines 11536--11536. Prerequisite pointer invoked at proof line(s) 19014; it names a standard fact of a sibling chapter that the counted argument uses, and every step of the counted proof built on it is retained verbatim in the delivered body.\par\medskip
\begin{lemma}
\label{lemma-projection-formula-proper}
Let $f : X \to Y$ be a proper morphism of schemes.
Let $E \in D^+(X_\etale)$ have torsion cohomology sheaves.
Let $K \in D^+(Y_\etale)$. Then
$$
Rf_*E \otimes_\mathbf{Z}^\mathbf{L} K =
Rf_*(E \otimes_\mathbf{Z}^\mathbf{L} f^{-1}K)
$$
in $D^+(Y_\etale)$.
\end{lemma}

\begin{proof}
There is a canonical map from left to right by
Cohomology on Sites, Section \ref{sites-cohomology-section-projection-formula}.
We will check the equality on stalks.
Recall that computing derived tensor products commutes with pullbacks.
See Cohomology on Sites, Lemma
\ref{sites-cohomology-lemma-pullback-tensor-product}.
Thus we have
$$
(E \otimes_\mathbf{Z}^\mathbf{L} f^{-1}K)_{\overline{x}}
=
E_{\overline{x}} \otimes_\mathbf{Z}^\mathbf{L} K_{\overline{y}}
$$
where $\overline{y}$ is the image of $\overline{x}$ in $Y$.
Since $\mathbf{Z}$ has global dimension $1$ we see that
this complex has vanishing cohomology in degree $< - 1 + a + b$
if $H^i(E) = 0$ for $i \geq a$ and $H^j(K) = 0$ for $j \geq b$.
Moreover, since $H^i(E)$ is a torsion abelian
sheaf for each $i$, the same is true for the cohomology
sheaves of the complex $E \otimes_\mathbf{Z}^\mathbf{L} K$.
Namely, we have
$$
(E \otimes_\mathbf{Z}^\mathbf{L} f^{-1}K)
\otimes_{\mathbf{Z}}^\mathbf{L} \mathbf{Q} =
(E \otimes_\mathbf{Z}^\mathbf{L} \mathbf{Q})
\otimes_{\mathbf{Q}}^\mathbf{L}
(f^{-1}K \otimes_{\mathbf{Z}}^\mathbf{L} \mathbf{Q})
$$
which is zero in the derived category.
In this way we see that Lemma \ref{lemma-proper-base-change-stalk}
applies to both sides to see that it suffices to show
$$
R\Gamma(X_{\overline{y}},
E|_{X_{\overline{y}}} \otimes_\mathbf{Z}^\mathbf{L}
(X_{\overline{y}} \to \overline{y})^{-1}K_{\overline{y}}) =
R\Gamma(X_{\overline{y}},
E|_{X_{\overline{y}}}) \otimes_\mathbf{Z}^\mathbf{L} K_{\overline{y}}
$$
This is shown in Lemma \ref{lemma-pull-out-constant}.
\end{proof}

\end{document}
